\begingroup
\setlength{\abovedisplayskip}{6pt}
\setlength{\belowdisplayskip}{6pt}
\setlength{\abovedisplayshortskip}{4pt}
\setlength{\belowdisplayshortskip}{4pt}
\setlength{\intextsep}{6pt}

\textbf{1. Input data.} The specified sound speed and air density are
\[
\boxed{c=340\,\mathrm{m/s},\qquad \rho=1.2\,\mathrm{kg/m^3}.}
\]
The specified pressure, loudspeaker diameter, and sound levels are
\[
\boxed{p=88\,\mathrm{mPa}=0.088\,\mathrm{Pa},\qquad
D=1.25\,\mathrm{m},\qquad L_{\mathrm{sp}}=72\,\mathrm{dB},\qquad
L_{\mathrm{horn}}=124\,\mathrm{dB}.}
\]

\textbf{2. Acoustic impedance.} The plane-wave acoustic impedance is
\[
Z=\rho c=(1.2)(340)=\boxed{408\,\mathrm{rayl}.}
\]

\textbf{3. Sound intensity for $p=88\,\mathrm{mPa}$.} Following the course-note convention,
\[
I=\frac{p^2}{Z}=\frac{(0.088)^2}{408}
  =\boxed{1.898\times10^{-5}\,\mathrm{W/m^2}}.
\]

\textbf{4. Reference intensity and pressure.} The reference intensity is
\[
\boxed{I_0=10^{-12}\,\mathrm{W/m^2}=1\,\mathrm{pW/m^2}.}
\]
Using the specified $Z$ without rounding,
\[
p_0=\sqrt{ZI_0}=\sqrt{(408)(10^{-12})}
 \approx\boxed{20.2\,\mu\mathrm{Pa}}.
\]

\textbf{5. Decibel values for the specified pressure.} The intensity and pressure levels are
\[
L_I=10\log_{10}\left(\frac{I}{I_0}\right)=72.8\,\mathrm{dB},
\qquad
L_p=20\log_{10}\left(\frac{p}{p_0}\right)=72.8\,\mathrm{dB}.
\]
Therefore,
\[
\boxed{L_I=L_p=72.8\,\mathrm{dB}.}
\]

\textbf{6. Equality of the two decibel values.} Since $I=p^2/Z$ and $I_0=p_0^2/Z$,
\[
\frac{I}{I_0}=\left(\frac{p}{p_0}\right)^2,
\qquad
10\log_{10}\left(\frac{I}{I_0}\right)=20\log_{10}\left(\frac{p}{p_0}\right).
\]
The two levels therefore agree when the reference pressure satisfies $p_0^2=ZI_0$.

\textbf{7. Loudspeaker at $72\,\mathrm{dB}$.} The intensity corresponding to the loudspeaker level is
\[
I_{\mathrm{sp}}=I_0 10^{L_{\mathrm{sp}}/10}
 =\boxed{15.85\,\mu\mathrm{W/m^2}}.
\]
The corresponding pressure and circular-face area are
\[
p_{\mathrm{sp}}=\sqrt{ZI_{\mathrm{sp}}}=\boxed{80.4\,\mathrm{mPa}},
\qquad
A=\frac{\pi D^2}{4}=1.23\,\mathrm{m^2}.
\]
Assuming the stated intensity is uniform over that area, the sound power is
\[
P_{\mathrm{sp}}=I_{\mathrm{sp}}A=\boxed{19.45\,\mu\mathrm{W}}.
\]

\textbf{8. Pressure from a $124\,\mathrm{dB}$ car horn.} The horn pressure is
\[
p_{\mathrm{horn}}=p_0 10^{L_{\mathrm{horn}}/20}
 =\boxed{32.0\,\mathrm{Pa}}.
\]
For comparison, the corresponding intensity is $I_{\mathrm{horn}}=2.51\,\mathrm{W/m^2}$.

\textbf{9. Pain-threshold comparison.} The course-note pain references are
\[
\boxed{I_{\mathrm{pain}}=1\,\mathrm{W/m^2},\qquad
p_{\mathrm{pain}}=\sqrt{ZI_{\mathrm{pain}}}=20.2\,\mathrm{Pa}.}
\]
The pain level is $120\,\mathrm{dB}$. The loudspeaker is $48\,\mathrm{dB}$ below this level, with intensity ratio $1.585\times10^{-5}$ and pressure ratio $3.981\times10^{-3}$. The horn is $4\,\mathrm{dB}$ above the pain level, with intensity ratio $2.51$ and pressure ratio $1.585$; its pressure is therefore $32.0\,\mathrm{Pa}$, above the $20.2\,\mathrm{Pa}$ pain reference.

\textbf{10. Graduate work.} Intensity is proportional to pressure squared at fixed acoustic impedance. The intensity level therefore uses $10\log_{10}$, while the pressure level uses $20\log_{10}$,
\[
10\log_{10}\left(\frac{I}{I_0}\right)
=20\log_{10}\left(\frac{p}{p_0}\right).
\]

\clearpage
\matlabcaption{MATLAB calculations for B.6.}
\lstinputlisting[style=hwmatlab]{matlab/B6.m}
\endgroup
